% ------------------------------------------------------------------
%  Clustering method: DBSCAN.
%  Written from template/method_template.tex; the part order is fixed.
%  Code counterpart:
%    xxcluster/cluster/partitional/density_based/dbscan.py
% ------------------------------------------------------------------
\subsubsection{DBSCAN}
\label{sec:tech:dbscan}

\paragraph{Overview}
\label{sec:tech:dbscan:overview}
DBSCAN defines a cluster as a maximal set of points connected through
regions of sufficient density, and belongs to the density-based family of
\textbf{Sect.~\ref{sec:background:families}}~\cite{ref_32}. Two properties
distinguish it from every other method reported here. The number of
clusters is a result rather than a request, so it is not asked to commit
to $|\mathcal C|$ in advance~\cite{ref_1}. And it may decline to assign an
observation at all, labelling it noise. That refusal is the reason it is
on this list: a method that assigns everything must fold a handful of
upset conditions into some regime, and the resulting scenario would
describe an operating state the plant never held.

\paragraph{Assumptions and applicability}
\label{sec:tech:dbscan:assumptions}
The assumption is local density, not shape: clusters may be elongated,
curved or nested, which is what a Voronoi partition cannot express. What
it does assume is that one density threshold separates cluster from noise
across the whole dataset. Where genuine regimes differ in density, a
single $(\varepsilon, \mathit{minPts})$ pair cannot recover both --- it
either splits the sparse one or merges the dense ones~\cite{ref_33}, and
that limitation is recorded in
\textbf{Sect.~\ref{sec:methodology:threats}}. The neighbourhood is a
distance in feature units, so the method is not scale-invariant and the
scaling step of \textbf{Sect.~\ref{sec:data:preprocessing}} decides which
operating variables drive the result.

\paragraph{Formulation}
\label{sec:tech:dbscan:formulation}
Let $N_\varepsilon(\mathbf p) = \{\mathbf q : d(\mathbf p, \mathbf q) \le
\varepsilon\}$ be the neighbourhood of $\mathbf p$, itself included. Three
definitions follow~\cite{ref_32}:
\begin{description}[itemsep=2pt,topsep=2pt]
    \item[Core point] $\mathbf p$ with $|N_\varepsilon(\mathbf p)| \ge
    \mathit{minPts}$.
    \item[Directly density-reachable] $\mathbf q$ from $\mathbf p$ when
    $\mathbf q \in N_\varepsilon(\mathbf p)$ and $\mathbf p$ is a core
    point. The requirement on $\mathbf p$ alone is what makes the relation
    asymmetric.
    \item[Density-connected] $\mathbf p$ and $\mathbf q$ when some core
    point reaches both through a chain of directly density-reachable
    steps.
\end{description}
A cluster is a maximal density-connected set. A non-core point inside the
radius of a core point joins it as a \emph{border point} but does not
extend the chain, which is what stops two clusters bridged by sparse
points from merging. An observation in no such set is noise, labelled
$-1$, which is the package-wide convention for an unassigned observation.

\paragraph{Hyperparameters and tuning}
\label{sec:tech:dbscan:hyperparameters}
Two, and neither has a default worth trusting. $\mathit{minPts}$ is the
smallest group the analyst is willing to call a regime rather than a rare
state, so it is a modelling choice before it is a parameter;
\cite{ref_33} recommends at least $n + 1$ for $n$ features and larger for
noisy data. $\varepsilon$ is a distance in the units of the scaled
features, chosen from the $k$-distance curve for $k = \mathit{minPts}$,
whose knee marks where points stop having close neighbours. Both are swept
rather than tuned to a favourable answer, and the sweep is over the
density parameters, not over $|\mathcal C|$: the selector of
\textbf{Sect.~\ref{sec:methodology:k}} does not apply to this family by
construction, since $|\mathcal C|$ is not an input it could vary.

\paragraph{Complexity}
\label{sec:tech:dbscan:complexity}
The cost is dominated by the radius query. With a spatial index over
low-dimensional data this is $\mathcal{O}(m \log m)$ time; without one, or
where the index degrades with dimension, it falls back to the
$\mathcal{O}(m^2)$ of computing every pair~\cite{ref_33}. Space is
$\mathcal{O}(mk)$ for the neighbourhoods, $k$ the mean number of
neighbours within $\varepsilon$, which grows quickly if $\varepsilon$ is
set too generously. A precomputed dissimilarity matrix is
$\mathcal{O}(m^2)$ in space by construction, so it is the route for a
custom measure rather than for volume.

\paragraph{Application to \projectname{} data}
\label{sec:tech:dbscan:application}
Implemented natively against the definitions above rather than adapted,
since they are the content of this section; the only borrowed component is
the radius neighbourhood query, which is an index and not the algorithm.
The implementation sets \texttt{labels\_} and the family base recounts
$|\mathcal C|$ and the noise from them, so the two cannot disagree with
the partition. Agreement with the reference implementation is asserted in
the test suite, labels and core points alike, rather than assumed.
Noise observations are reported with their position, their timestamp where
the dataset carries one, and how many neighbours short of the threshold
they fell --- a point one neighbour short is a different finding from one
entirely alone. Whether any of them becomes a rare-event scenario is not
decided here; that belongs to the scenario task, and must be recorded
there.

\paragraph{Results}
\label{sec:tech:dbscan:results}
\placeholder{Every index of Section~\ref{sec:methodology:metrics} over the
swept $(\varepsilon, \mathit{minPts})$ grid, the noise fraction at each
setting, and the timestamps of the noise observations at the chosen one.}

\paragraph{Limitations}
\label{sec:tech:dbscan:limitations}
One global density threshold, as above: regimes of differing density
cannot both be recovered by one parameter pair, which is the case
hierarchical density methods exist to answer~\cite{ref_1}. The result is
sensitive to $\varepsilon$, and the sensitivity is not symmetric --- a
slightly generous radius merges regimes, a slightly tight one converts
their edges to noise. A border point in reach of two clusters joins
whichever is scanned first, so its label is an artefact of order rather
than of the data; the count of such points is worth reporting where the
partition is close. Finally, the method is transductive: it labels the
sample it was given and offers no rule for an unseen observation, so it
cannot score a held-out period without refitting.

\paragraph{Summary}
\label{sec:tech:dbscan:summary}
The only method reported here that can say an operating state is too rare
to be a regime, and the only one whose cluster count is discovered rather
than requested. It earns its place on the scenario path by that refusal
rather than by scoring well on indices that assume the compact geometry it
is meant to escape --- a caveat carried into
\textbf{Sect.~\ref{sec:comparison}}.

% ==================================================================
%  REFERENCES USED
%
%    ref_1    Singh & Singh (2024), density-based section
%                                cited in: Overview, Limitations
%    ref_32   Ester, Kriegel, Sander & Xu (1996), A density-based
%             algorithm for discovering clusters in large spatial
%             databases with noise
%                                cited in: Overview, Formulation
%    ref_33   Schubert et al. (2017), DBSCAN revisited, revisited
%                                cited in: Assumptions, Hyperparameters,
%                                          Complexity
%
%  Read, not cited:
%    ref_18   Kapanski et al. (2025), density clustering for temporal
%      segmentation of urban water consumption -- the precedent for
%      reporting noise with its period rather than discarding it. To be
%      cited in Results once the run on project data is reported.
% ==================================================================
